\documentclass[11pt]{article}
\usepackage[margin=1in,top=0.85in,bottom=0.8in]{geometry}
\usepackage{amsmath,amssymb}
\usepackage[protrusion=true,expansion=false]{microtype}
\usepackage{xcolor}
\usepackage{enumitem}
\definecolor{acc}{HTML}{1B6B58}
\definecolor{grey}{HTML}{5C6763}
\newcommand{\R}{\mathbb{R}}
\newcommand{\E}{\mathbb{E}}
\newcommand{\Cov}{\operatorname{Cov}}
\newcommand{\Var}{\operatorname{Var}}
\newcommand{\Nor}{\mathcal{N}}
\newcommand{\tr}{^{\!\top}}
\pagestyle{empty}
\setlength{\parindent}{0pt}
\setlength{\parskip}{0.4em}

\begin{document}

{\footnotesize\color{acc}\textsc{Probability for Modern Machine Learning \quad$\cdot$\quad Quiz 1 \quad$\cdot$\quad Wednesday, September 9}}

\vspace{0.4em}
{\large\bfseries Quiz 1}\hfill {\color{grey}15 minutes \quad$\cdot$\quad closed book \quad$\cdot$\quad no calculators}

\vspace{0.5em}
\hrule height 0.5pt
\vspace{0.3em}

{\footnotesize\color{grey} Three questions, five minutes each. Drawn from the core problems of Homeworks 1 and 2. Answers should be short; a correct one-line argument is worth more than a page of work.}

\vspace{1.4em}

\textbf{1.} \ \textsc{\small compute} \quad Let $X\sim\Nor(0,\sigma^2)$ on $\R$.

\begin{enumerate}[label=(\alph*),leftmargin=1.8em,itemsep=0.35em,topsep=0.35em]
\item Write down the differential entropy $h(X)$.
\item A colleague reports that a certain real-valued random variable has differential entropy $-3$. Is this possible? Give an example or explain why not.
\end{enumerate}

\vspace{1.8em}

\textbf{2.} \ \textsc{\small compute} \quad Observations follow $y=Hw+\sigma\xi$, with $\xi\sim\Nor(0,I_n)$ and a prior $w\sim\Nor(0,\tau^2I_p)$ independent of $\xi$.

\begin{enumerate}[label=(\alph*),leftmargin=1.8em,itemsep=0.35em,topsep=0.35em]
\item Write down $\E[w\mid y]$.
\item Identify the regularization parameter in terms of $\sigma$ and $\tau$, and say in one sentence what a large value of it means.
\end{enumerate}

\vspace{1.8em}

\textbf{3.} \ \textsc{\small which hypothesis fails} \quad Let $g\sim\Nor(0,1)$ and let $\varepsilon$ be an independent random sign. Set
\[
X=g,\qquad Y=\varepsilon g .
\]

\begin{enumerate}[label=(\alph*),leftmargin=1.8em,itemsep=0.35em,topsep=0.35em]
\item Show $\Cov(X,Y)=0$.
\item Are $X$ and $Y$ independent? Justify in one sentence.
\item Both $X$ and $Y$ are standard normal and they are uncorrelated. Name precisely the hypothesis of the theorem from Lecture 1 that fails here.
\end{enumerate}

\vspace{2.2em}
\hrule height 0.4pt
\vspace{0.3em}
{\footnotesize\color{grey} Name:\hspace{16em}\hrulefill}

\end{document}
